\section{Riemannian comparison at a prescribed conformal scale}
\label{r:section}

The comparison uses the length metric $\mathfrak h_r=X_r^2h_r$
on the meridian of a polarized metric $g_r=h_r+X_r^2\dd\phi^2$.
Thus rotation of a meridional arc of $\mathfrak h_r$-length $\ell$
has area $2\pi\ell$.  The subscript $r$ labels the comparison metric;
it never denotes differentiation.  The potentials, their normalization,
and the target metric $G$ are those of the axial reduction.  In particular,
\[
 \Phi=(u,\psi,\chi,z),\qquad u=\log X_r,
 \qquad G=\dd u^2+|\cD_{X_r}|^2.
\]

\begin{proposition}[Riemannian comparison]\label{r:comparison}
Let $\Omega_r$ be the closure of the component toward infinity obtained
by removing finitely many smooth compact invariant obstacles from a
rotation exterior as in the axial reduction.  The omitted region includes
the original inner side, and $\partial\Omega_r$ is nonempty.  Suppose the
following hold.
\begin{enumerate}[label=\textup{(\roman*)}]
\item The polarized metric $g_r$ is smooth up to its inner faces,
asymptotically flat of order $q_d>1/2$ through at least two derivatives,
and has finite ADM energy $E_r$.
\item The potentials are restrictions of potentials on the original
rotation exterior.  Their constants on its two axis sides are
$(-Q,0,-2J)$ and $(Q,0,2J)$, and they have the polar and end regularity
of Lemma~\ref{a:axis-regularity} and Equation~\eqref{a:end-regularity}.
\item One has
\begin{equation}\label{r:scalar-hypothesis}
 R_{g_r}\ge 2|\cD_{X_r}|_{g_r}^2,\qquad R_{g_r}>0.
\end{equation}
Each inner face has strictly negative mean curvature for the normal
toward infinity.  Alternatively, the inner boundary is already a smooth
invariant stable minimal boundary.
\item Every meridional path in the retained region joining the two
original axis sides has $\mathfrak h_r$-length at least $L>0$.
The same assumption may be verified by a length-decreasing projection
from any added collars to a reference meridian with that lower bound.
\end{enumerate}
For every $R$ on the open physical branch
$R^2>\sqrt{Q^4+4J^2}$, let $M=F_*(R)$ and $b>0$ be the model
parameters in Lemma~\ref{a:model-parameters}.  Then
\begin{equation}\label{r:mass-bound}
 E_r\ge F_*(R)+\frac b2\log\frac{L}{2R^2}.
\end{equation}
If the given inner boundary is a single stable minimal sphere, the strict
inequality $R_{g_r}>0$ may be omitted.
\end{proposition}

Here and below, meridional paths can be taken piecewise smooth, with
ordinary axis endpoints and compact image away from the punctures of
any cylindrical completion.  The smooth approximation argument below
also gives the same lower bound for their ordinarily rectifiable limits.

\subsection{Stable boundary and preservation of crossing lengths}
\label{r:hulls}

We first construct a stable minimal boundary when the initial faces have
strictly negative mean curvature.  Choose a sufficiently distant coordinate
sphere in a strictly mean-convex foliation of the end.  In the bounded
region between this sphere and the inner obstacles, minimize the full
perimeter of filled inner sets containing the obstacles.  A smooth fill
on the excluded side may be used to pose this ambient perimeter problem;
it asserts nothing about the constraints on that side.  Perimeter
compactness and lower semicontinuity give a minimizer.  Smooth-obstacle
regularity in dimension three gives a $C^{1,1}$ boundary, smooth and
minimal away from the obstacle
\cite[Regularity Theorem~1.3(iii)]{HuiskenIlmanen2001}.

Neither barrier can be touched.  At an inner contact point, write both
surfaces as graphs with the common normal toward infinity.  The minimizing
graph lies above the obstacle graph and has weak outward mean curvature
at least zero, by outward variations, whereas the obstacle has strictly
negative mean curvature.  The graph comparison principle excludes this
contact.  The opposite strict comparison excludes contact with the outer
sphere.  The mean-convex foliation also prevents escape beyond the chosen
outer barrier.  The free boundary is therefore a smooth stable minimal
surface.

We can choose the minimizer invariantly.  Among the minimizers choose one
with largest filled volume; such a choice exists by compactness.  If
$E_1,E_2$ are two minimizers, perimeter submodularity gives
\[
 \operatorname{Per}(E_1\cup E_2)+\operatorname{Per}(E_1\cap E_2)
 \le \operatorname{Per}(E_1)+\operatorname{Per}(E_2).
\]
Both union and intersection are admissible, so both are minimizers.
The union of two distinct largest-volume minimizers would have larger
volume.  Consequently that minimizer is unique up to null sets, and
invariance of the problem under rotations makes it invariant.  Fill
bounded complementary pockets and retain the frontier toward infinity.
If a stable minimal boundary was supplied initially, begin with that
frontier instead.

For a two-sided stable minimal component $\Sigma$, testing its Jacobi
form by $1$ and using the Gauss equation gives
\begin{equation}\label{r:stable-sphere}
 \int_\Sigma R_\Sigma\,\dd A
 \ge \int_\Sigma(R_{g_r}+|\mathrm{II}|^2)\,\dd A>0.
\end{equation}
Thus every orientable component is a sphere.  Strict scalar positivity is
used only here, which explains the last assertion of the proposition.
The rotation action on each sphere has two poles.  In the capped rotation
topology of the axial reduction, the corresponding meridional arcs and
the finite axis intervals between their poles bound their compact-side
disks.  The filled frontier components are not nested.  Exactly one
compact side contains the cap behind the original boundary.  Its arc
joins the two original axis sides; the endpoints of every other arc lie
on a single original side.  Since the potentials are inherited from the
whole original exterior, all three potential jumps on these additional
spheres vanish.

For later use, we spell out the conversion of crossing lengths to
enclosing areas.  A piecewise smooth meridional path joining the original
axis sides can first be perturbed off the inner boundary and off the axis
except at its endpoints, with arbitrarily small increase of length in
the continuous metric $X_r^2h_r$.  Arrange finitely many transverse
self-intersections, discard loops, and extract a simple crossing arc.
At its endpoints round it to meet the axis orthogonally in ordinary polar
coordinates.  The change in weighted length tends to zero because $X_r$
vanishes linearly there.  Rotation produces a smooth invariant enclosing
cut.  In particular, if every invariant enclosing cut in a reference
exterior has area at least $A_0$, then every crossing path there has length
at least
\begin{equation}\label{r:area-length}
 L=\frac{A_0}{2\pi}.
\end{equation}
This argument uses invariant cuts only.  A lower bound for all enclosing
cuts supplies the required bound for that subclass.

\subsection{Cylindrical completion and its curvature}
\label{r:cylinders}

Discard the compact side of the stable boundary.  Warped-product and
two-dimensional conformal curvature formulas give, off the axis,
\[
 R_{g_r}=2K_{h_r}-2X_r^{-1}\Delta_{h_r}X_r,
 \qquad
 K_{\mathfrak h_r}=X_r^{-2}(K_{h_r}-\Delta_{h_r}u).
\]
Since $X_r^{-1}\Delta_{h_r}X_r=\Delta_{h_r}u+|\dd u|_{h_r}^2$,
we obtain
\begin{equation}\label{r:curvature}
 X_r^2K_{\mathfrak h_r}
 =\frac12R_{g_r}+|\dd u|_{h_r}^2,
 \qquad
 K_{\mathfrak h_r}\ge |\dd\Phi|_{\mathfrak h_r,G}^2.
\end{equation}
Minimality of a rotated surface is precisely the geodesic equation for
its generating arc in $\mathfrak h_r$, since its area is $2\pi$ times
that arc's length.

Fix a stable boundary sphere and let $\ell$ be arclength along its
generating arc in $\mathfrak h_r$.  Let $p>0$ be its invariant first
Jacobi eigenfunction, with eigenvalue $\lambda\ge0$, and put $j=X_rp$.
The eigenfunction is invariant because the positive first eigenfunction
is unique up to scale.  A sphere normal speed $s_0$ induces geodesic
normal speed $X_rs_0$.  The corresponding second-variation bilinear
forms therefore satisfy
\[
 I_\Sigma(p,s_0)=2\pi I_{\mathrm{geo}}(X_rp,X_rs_0).
\]
Since the rotational area measure is $2\pi\,\dd\ell$, testing by
invariant $s_0$ supported away from the poles gives
\[
 \int(-j''-K_{\mathfrak h_r}j)X_rs_0\,\dd\ell
 =\lambda\int ps_0\,\dd\ell.
\]
Consequently, on the open generating arc,
\begin{equation}\label{r:jacobi-cylinder}
 -j''-K_{\mathfrak h_r}j=\lambda\frac{j}{X_r^2}\ge0.
\end{equation}

Attach the half-cylinder $\tau\ge0$ with length metric
\begin{equation}\label{r:cylinder-metric}
 \dd\ell^2+j(\ell)^2\dd\tau^2,
\end{equation}
and hold $\Phi$ constant along each $\tau$-line.  Its Gauss curvature is
$-j''/j$.  Equations~\eqref{r:curvature} and
\eqref{r:jacobi-cylinder} show that this curvature dominates the original
full map energy on the seam, and hence dominates the tangential map
energy of the extended map.  Thus the second inequality in
\eqref{r:curvature} holds on the cylinder too.

The seam is geodesic on both sides.  Match normal collars so that metric
coefficients and the map are continuous there.  The resulting coefficients
are Lipschitz and piecewise smooth.  No negative curvature measure occurs
at the seam: the distributional curvature includes the sum of the two
geodesic-curvature boundary terms, which is zero.  More explicitly, away
from the axis choose a continuous piecewise smooth oriented orthonormal
frame whose first vector is tangent to the seam.  The pullback of its
connection form to the seam vanishes on each side by geodesicity.
Stokes' formula therefore has no seam term.  A conformal-coordinate frame
differs from this frame by a continuous $W^{1,p}$ angle; its connection
form changes by the differential of that angle, so the same distributional
curvature identity holds in conformal coordinates.

The apparent degeneracy at a pole is only the usual axial degeneracy.
If $s_B$ is ordinary meridional arclength on the sphere, then
$\dd\ell=X_r\dd s_B$, and the nonsingular meridional metric underlying
\eqref{r:cylinder-metric} is
\begin{equation}\label{r:regular-collar}
 \dd s_B^2+p(s_B)^2\dd\tau^2.
\end{equation}
Here $p$ is smooth and even at each pole, and $X_r$ is smooth and odd in
the reflected polar coordinate, with unit first derivative.  Thus the
three-dimensional lift is regular at its axes.  Extend the original
orbit arc across each pole by reflection and match normal collars of the
two nonsingular metrics.  The matched metric remains continuous and
piecewise smooth, preserves reflection parity, and satisfies the no-cone
condition.  This construction also supplies reflected Lipschitz charts
where a seam meets the axis.

Projection $(\ell,\tau)\mapsto(\ell,0)$ decreases length in
\eqref{r:cylinder-metric}.  Hence every excursion into an added cylinder
can be replaced by an arc of its attaching boundary.  Its projection is
rectifiable also in the ordinary base metric, including at the poles,
so the smoothing argument preceding \eqref{r:area-length} applies.
After conformal filling, the end of the sphere separating the original
obstacle becomes one distinguished puncture.  The two sides between that
puncture and infinity are exactly the original two axis sides.  Every
additional puncture connects adjacent intervals on a single side.
It follows that all crossing paths in the completed meridian still
have length at least $L$.

\subsection{Uniformizing the completed meridian}
\label{r:uniformization}

The conformal structure is that of the nonsingular meridional metric
$h_r$, equivalently that of $\mathfrak h_r$ off the axis.  Reflect it
across all ordinary axis intervals.  Each product end has an explicit
conformal puncture coordinate.  Indeed, with
\[
 t_B=\int^{s_B}\frac{\dd s}{p(s)},
 \qquad T_B=\int_{\text{pole}}^{\text{pole}}\frac{\dd s}{p(s)}\in(0,\infty),
\]
Equation~\eqref{r:regular-collar} is
$p^2(\dd t_B^2+\dd\tau^2)$.  Exponentiating a strip coordinate with
angular variable $\alpha=\pi t_B/T_B$ fills the cylindrical end by a
boundary point.  In logarithmic coordinates the length factor is
$(T_B/\pi)X_rp$, independent of the longitudinal coordinate.  Dividing
it by $\sin\alpha$ gives a positive regular function through both
endpoints.  This verifies the finite-width cylindrical profiles required
by Lemma~\ref{c:cut}.

At infinity, asymptotic flatness conformally compactifies the reflected
end after isothermal rectification.  The compactified reflected surface
is a sphere with finitely many points removed before filling; its genus
is zero by the capped meridional topology.  Uniformization, chosen
compatibly with reflection, gives an upper half-plane coordinate $Z$.
The reflection becomes complex conjugation.  Put the distinguished
cylindrical puncture at $0$, infinity at infinity, and normalize the
Euclidean scale at infinity to one.  Other cylindrical ends are marked
points on the real boundary.

We record the regularity and mass normalization of these coordinates.
On compact sets, including reflected seams, the conformal coefficients
are uniformly elliptic and Lipschitz.  Local Beltrami estimates give
isothermal changes in $W^{2,p}_{\mathrm{loc}}$ for every finite $p$,
hence $C^{1,\lambda}_{\mathrm{loc}}$ for every $\lambda<1$.
Their first derivatives do not vanish.  This regularity justifies the
connection-form calculation above and gives the weak conformal curvature
formula across a seam.

For any
\begin{equation}\label{r:decay-exponent}
 \frac12<q'<\min(q_d,1),
\end{equation}
the normalized end transition from the original Euclidean meridian has
the expansion
\begin{equation}\label{r:af-isothermal}
 Z=Z_{\mathrm{Eucl}}+O_2(\rho_\infty^{1-q'}),
\end{equation}
after a Euclidean motion of the original chart.
To see the claimed rate, invert in the original end chart.  The reflected
Beltrami coefficient is $O(|Z|^{q_d})$ at the inverted origin, and its
derivative estimates make it $C^{0,q'}$ after extension there by zero.
Local rectification has a nonzero conformal first derivative and a
$C^{1,q'}$ expansion at that origin.  Inverting back gives the asserted
zeroth- and first-order asymptotics.  On rescaled original end annuli,
first-order elliptic Schauder estimates for the Beltrami equation give
the required second derivatives of the error, decreasing $q'$ if needed.
Reflection preserves the even longitudinal and odd transverse estimates.
Lifting these changes as a longitudinal change and a polar radial change
therefore gives admissible three-dimensional asymptotically Euclidean
coordinates of order $q'>1/2$.

These coordinates have the same ADM energy.  Here is the flux argument
at exactly the regularity in use.  If $h=g_r-\delta$ in the old end chart,
scalar-curvature linearization gives
\[
 R_{g_r}=\partial_i\partial_j h_{ij}-\Delta h_{ii}
                  +O(\rho_\infty^{-2q_d-2}).
\]
The error is integrable, and finite mass flux together with $R_{g_r}\ge0$
implies integrability of $R_{g_r}$ and of the Euclidean scalar
linearization.  Under a change by
$O_2(\rho_\infty^{1-q'})$, the leading change in metric perturbation is
a Euclidean symmetric gradient, plus the old perturbation in the old
coordinates.  The linear mass flux of a symmetric gradient is zero:
extend its vector field smoothly inside a sphere and commute Euclidean
derivatives in its scalar linearization.  All quadratic flux errors are
$O(\rho_\infty^{1-2q'})=o(1)$.  Surface and derivative changes in the
remaining old perturbation have the same vanishing bound.  Finally the
integrable scalar linearization identifies its flux on the resulting
asymptotically round deformed spheres with its flux on coordinate spheres.
Thus $E_r$ is unchanged.

Ordinary axes have a simple linear zero of the length density; the
cylindrical profiles were just checked; the outer density has the
Euclidean quadratic growth.  All assumptions of Lemma~\ref{c:cut}
therefore hold.  Fix a model $R$ and apply that lemma with capacity
$k=b/2$ and with $0<L'<L$.  Write its coordinates as
\[
 Z=f(e^{\sigma+i\vartheta}),\qquad
 \sigma\ge0,\quad 0\le\vartheta\le\pi,
 \qquad f(\xi)=\frac b2\xi+O(1).
\]
The artificial inner cut avoids the additional punctures.  In these
coordinates put
\begin{equation}\label{r:log-density}
 \mathfrak h_r=e^{2q}(\dd\sigma^2+\dd\vartheta^2).
\end{equation}
The curvature inequality and the inner cut estimate become
\begin{equation}\label{r:strip-curvature-cut}
 -\Delta q\ge |\partial\Phi|_G^2,
 \qquad
 q(0,\vartheta)\ge\log\left(\frac{L'}2\sin\vartheta\right).
\end{equation}
The differential inequality is distributional across any retained seams.
A cut analytic in uniformizing coordinates need not be smooth across such
a seam in the original coordinates; the weak calculation below covers
this case.

\subsection{Target convexity and all finite boundary terms}
\label{r:weak-comparison}

Use the model map $\Phi_*$ of Proposition~\ref{a:model} on this
same strip, and set
\[
 H_0=q-q_*,\qquad w=\sinh\sigma\sin\vartheta.
\]
By Equation~\eqref{a:model-map}, the model satisfies
$-\Delta q_*=|\partial\Phi_*|_G^2$ and has
weighted tension zero,
$\sum_i\nabla_i^G(w\partial_i\Phi_*)=0$.
Because $(\R^4,G)$ is complete, simply connected and nonpositively
curved, there is a unique geodesic from $\Phi_*(x)$ to $\Phi(x)$,
depending smoothly on its endpoints.  Let $V(x)$ be its initial velocity
and define
\[
 C_i=\langle V,\partial_i\Phi_*\rangle_G.
\]
For its pointwise geodesic interpolation $\Phi_t$, nonpositive curvature
and the Jacobi-field equation imply convexity of
$|\partial\Phi_t|_G^2$.  The tangent-line inequality at $t=0$ gives
\[
 |\partial\Phi|_G^2-|\partial\Phi_*|_G^2
 \ge 2\sum_i\langle\nabla_i^G V,\partial_i\Phi_*\rangle_G
 =2w^{-1}\operatorname{div}(wC).
\]
The last equality is the weighted model equation.  By local Sobolev
approximation this remains valid weakly across the seams.  Subtracting
the two curvature equations yields
\begin{equation}\label{r:target-convexity}
 -\Delta H_0\ge2w^{-1}\operatorname{div}(wC).
\end{equation}
Since $\Delta w=0$, the vector field
\begin{equation}\label{r:comparison-flux}
 \mathcal F=-w\nabla H_0+H_0\nabla w-2wC
 \qquad\hbox{satisfies}\qquad \operatorname{div}\mathcal F\ge0.
\end{equation}
We now account for every finite boundary piece.

\paragraph{Ordinary axes.}
The no-cone condition in conformal meridional coordinates gives the trace
identity
\begin{equation}\label{r:axis-traces}
 H_0=2(u-u_*).
\end{equation}
Indeed the nonsingular conformal scale equals the transverse derivative
of $X_r$ at the axis, and the same statement holds for the model.
The common axis constants and polar vanishing orders of the potentials
give, in target norm,
\[
 V=(u-u_*)\partial_u+o(1),\qquad
 \sin\vartheta\,\partial_\vartheta\Phi_*
       =\cos\vartheta\,\partial_u+o(1).
\]
Thus the traces of $H_0\partial_\nu w$ and $-2wC_\nu$ cancel on both
axes, with their respective outward normals.  The remaining term
$w\partial_\nu H_0$ has zero limiting integral.

The last assertion also holds at axis--seam intersections.  Subtracting
the common $\log\sin\vartheta$ from the logarithmic densities leaves
continuous H\"older functions with locally $L^p$ gradients for finite
$p>1$.  To include the seam--axis intersection explicitly, take a common
signed transverse coordinate $y$ in the matched reflected collars.
The ratio $X_r/y$ is positive, continuous and piecewise smooth, with
bounded weak first derivatives: the values match on the seam, and
ordinary polar regularity and the no-cone condition give the same
continuous positive trace at the corner.  Write $\vartheta_{\rm ax}$
for $\vartheta$ at the lower axis and for $\pi-\vartheta$ at the upper
axis, extended by signed reflection.  In isothermal coordinates,
\[
 \frac{X_r}{\vartheta_{\rm ax}}
      =\frac{X_r}{y}\frac{y}{\vartheta_{\rm ax}}.
\]
The $W^{2,p}$ reflected coordinate change and the one-dimensional Hardy
estimate give $y/\vartheta_{\rm ax}\in W^{1,p}$; the ratio is bounded away from
zero.  Thus its logarithm and $\log(X_r/\vartheta_{\rm ax})$ are in $W^{1,p}$.
The nonsingular conformal factor is also in $W^{1,p}$, by pullback of
the Lipschitz metric under the same coordinate change.  This proves
the asserted regularity of the renormalized logarithmic density.
In an axis layer of
width $\delta$, $w\le C\delta$.  Paired with a cutoff derivative of size
$C/\delta$, the integral involving $w\nabla H_0$ is bounded by
$C\int_{\text{layer}}|\nabla H_0|=o(1)$.  For the two other terms,
continuity of their renormalized traces and \eqref{r:axis-traces}
give precisely the cancellation above.  Equivalently one may choose
offset axes along an exhaustion with the same vanishing integrated error.

\paragraph{Additional punctures.}
Every puncture still visible on a side of the strip has zero potential
jump.  Let $d$ denote its Euclidean distance in local flat coordinates.
The product cylinder and the conformal puncture coordinate give
\begin{equation}\label{r:puncture-estimates}
 |H_0|\le C(1+|\log d|),\qquad
 |\partial H_0|\le C/d,\qquad w\le Cd.
\end{equation}
There is also the uniform target-distance bound
$|V|_G\le C(1+|\log d|)$.  To check it without losing control near the
local polar directions, let $\beta$ be the local puncture angle.
The cylindrical radius $X_r$ is comparable to $\sin\beta$, while
$X_*$ is comparable to $d\sin\beta$.  At height $X_r$, first match the
two horizontal potential coordinates and then the central coordinate.
For clarity put
$\widetilde z=z+\psi_{\rm ax}\chi-\chi_{\rm ax}\psi-z_{\rm ax}$
and denote subtraction of model values by $\delta$.  The central
Heisenberg displacement is exactly
\begin{equation}\label{r:central-displacement}
 \delta z+\psi_*\delta\chi-\chi_*\delta\psi
 =\delta\widetilde z
  +(\psi_*-\psi_{\rm ax})\delta\chi
  -(\chi_*-\chi_{\rm ax})\delta\psi.
\end{equation}
The cylinder map is fixed on one compact smooth sphere.  Its horizontal
coordinates differ from the common axis values by $O(\sin^2\beta)$,
and its corrected central coordinate is $O(\sin^4\beta)$ uniformly.
The model has the corresponding orders in $d\sin\beta$.
Equation~\eqref{r:central-displacement} is therefore
$O(\sin^4\beta)$.  Division by $X_r$ for the horizontal coordinates
and by $X_r^2$ for the central coordinate gives bounded displacements.
This makes the target path a bounded-cost path uniformly in $\beta$.
Then change the height;
that costs at most $C+|\log d|$.  This proves the bound on $V$.
The model derivative satisfies $w|\partial\Phi_*|_G\le C$ near the
puncture, so
\[
 w|C|\le C(1+|\log d|).
\]
On a puncture semicircle of radius $d$, the three terms in
\eqref{r:comparison-flux} consequently have integrated magnitudes
$O(d)$, $O(d(1+|\log d|))$, and $O(d(1+|\log d|))$, respectively.
All tend to zero.  This is where the vanishing jumps of the additional
components are essential.

\paragraph{The artificial inner cut and its corners.}
On $\sigma=0$ the outward normal is $-\partial_\sigma$, and
$w=0$, $\partial_\nu w=-\sin\vartheta$.  The model normal derivative
$\partial_\sigma\Phi_*$ is bounded in target norm on finite radial
ranges, including the corners after the polar normalization.  The only
remaining inner flux is therefore
\[
 -\int_0^\pi H_0(0,\vartheta)\sin\vartheta\,\dd\vartheta.
\]
For a cut crossing seams, use the local $L^p$ gradient bounds and the
same width-$\delta$ argument as for the axes.  Products of the inner-edge
and axis cutoffs handle the corners, since $w$ vanishes on both edges.
This uses only the trace of the cut density; there is no normal-derivative
condition on the artificial cut.

For precision, first perform this cutoff integration in a finite strip
$0<\sigma<T$, with small disks about its finitely many marked side
punctures removed.  Send the nonpuncture edge-layer widths to zero, using
the preceding estimates.  Then shrink the puncture disks.  These operations
leave just the outer flux at $\sigma=T$ and the displayed inner flux.
With
\[
 \langle D\rangle=\frac12\int_0^\pi D\sin\vartheta\,\dd\vartheta,
\]
we obtain, for large regular $T$ and then in the limit,
\begin{equation}\label{r:integrated-comparison}
 \lim_{T\to\infty}
 \left\langle-\sinh T\,\partial_\sigma H_0
       +\cosh T\,H_0-2\sinh T\,C_\sigma\right\rangle
 \ge \langle H_0(0,\cdot)\rangle.
\end{equation}
The next calculation proves existence and identifies the value of the
outer limit.

\subsection{ADM flux and completion of the comparison}
\label{r:adm}

Set $\rho_0=(b/2)e^\sigma$.  Write the comparison metric near infinity as
\[
 h_r=e^{2P_1}(\dd\rho_0^2+\rho_0^2\dd\vartheta^2),\qquad
 X_r=\rho_0\sin\vartheta\,e^{P_2}.
\]
Both $P_1$ and $P_2$ are $O_1(\rho_0^{-q'})$ in the end sense,
with the reflected polar estimates.  A direct ADM calculation gives
\begin{equation}\label{r:radial-adm}
 E_r=\lim_{\rho_0\to\infty}\frac{\rho_0}{2}
 \left\langle-\partial_\sigma(P_1+P_2)+P_1-P_2\right\rangle.
\end{equation}
For clarity, in the Euclidean orthonormal spherical frame the metric
perturbation has entries $A,A,B$, with
$A=e^{2P_1}-1$ and $B=e^{2P_2}-1$.
The radial ADM integrand is exactly
\[
 -\partial_{\rho_0}(A+B)+(A-B)/\rho_0.
\]
Substitution of $A=2P_1+O(P_1^2)$ and $B=2P_2+O(P_2^2)$, followed by
integration over a sphere, gives \eqref{r:radial-adm}.  The discarded
flux is $O(\rho_0^{1-2q'})=o(1)$.

Denote subtraction of model values by $\delta$.  Equation
\eqref{r:log-density} gives
\begin{equation}\label{r:mass-differences}
 H_0=\delta(P_1+P_2),\qquad u-u_*=\delta P_2.
\end{equation}
We also claim
\begin{equation}\label{r:target-log-end}
 C_\sigma=u-u_*+o(\rho_0^{-1})
\end{equation}
uniformly in the polar angle.  To verify this estimate, normalize the
model point in the target by its Heisenberg translation and dilation,
both isometries of $G$.  In this fixed smooth target chart, the log-height
displacement is $O(\rho_0^{-q'})$, the normalized horizontal potential
displacement is $O(\rho_0^{-1})$, and the normalized central displacement
is $O(\rho_0^{-2})$.  The end estimates and common polar constants imply
these bounds also when $\sin\vartheta$ is small.  The horizontal
differences vanish to second polar order, and
Equations~\eqref{a:end-central} and~\eqref{r:central-displacement}
give the fourth-order bound for the translated central difference,
uniformly on rescaled end annuli.  These orders cancel the inverse
height factors.
The geodesic logarithm differs from these normalized coordinate
displacements by
$O(\rho_0^{-2q'}+\rho_0^{-2})$ in its radial component.  Finally,
\[
 \partial_\sigma\Phi_* =\partial_u+O(\rho_0^{-1})
\]
in the orthonormal target frame.  Pairing with $V$ proves
\eqref{r:target-log-end}, since $2q'>1$.

Since $\sinh\sigma$ and $\cosh\sigma$ equal
$\rho_0/b+O(\rho_0^{-1})$, Equations~\eqref{r:mass-differences}
and \eqref{r:target-log-end} turn the left side of
\eqref{r:integrated-comparison} into
\[
 \frac1b\lim_{\rho_0\to\infty}\rho_0
 \left\langle-\partial_\sigma\delta(P_1+P_2)
                    +\delta P_1-\delta P_2\right\rangle
 =\frac{2(E_r-M)}b.
\]
At the inner edge the model identity is
$q_*(0,\vartheta)=\log(R^2\sin\vartheta)$.  By
\eqref{r:strip-curvature-cut},
\[
 H_0(0,\vartheta)\ge\log\frac{L'}{2R^2}.
\]
Thus $E_r\ge M+(b/2)\log(L'/(2R^2))$.  Keeping the model $R$ fixed
and letting $L'\uparrow L$ proves Proposition~\ref{r:comparison}.

\begin{corollary}[Optimization on the physical branch]\label{r:optimized}
Under the hypotheses of Proposition~\ref{r:comparison}, if
$L/2\ge\sqrt{Q^4+4J^2}$, then
\begin{equation}\label{r:optimized-bound}
 E_r\ge F_*\!\left(\sqrt{L/2}\right).
\end{equation}
In particular, with $L=A_0/(2\pi)$ this becomes
\[
 E_r\ge\sqrt{\frac{A_0}{16\pi}+\frac{Q^2}{2}
                  +\frac{\pi(Q^4+4J^2)}{A_0}}.
\]
\end{corollary}

\begin{proof}
Put $C=Q^4+4J^2$.  The model parameter relations give
\[
 M=\frac12\sqrt{R^2+2Q^2+C/R^2},\qquad
 b=\frac{R^4-C}{4MR^2},\qquad M'(R)=\frac bR.
\]
Consequently
\begin{equation}\label{r:optimization-derivative}
 \frac{\dd}{\dd R}\left(M+\frac b2\log\frac{L}{2R^2}\right)
 =\frac{b'(R)}2\log\frac{L}{2R^2}.
\end{equation}
Here $b'(R)>0$ by Lemma~\ref{a:model-parameters}.
If $L/2>\sqrt C$, Equation~\eqref{r:optimization-derivative} shows
that the comparison expression increases up to $R^2=L/2$ and decreases
afterward.  At that point its logarithmic term vanishes, proving
\eqref{r:optimized-bound}.  If $L/2=\sqrt C$, use models with
$R^2\downarrow L/2$; then $b\downarrow0$ and the logarithmic correction
tends to zero.  The limiting argument is numerical and uses only
strictly subextremal comparison models.
\end{proof}
